% paper6_methods.tex -- Materials and Methods of Paper 6, input by
% papers/paper6_trials.tex.
\section{Materials and Methods}

\subsection{Sites}
The trials were laid out at three sites chosen to represent contrasting
climatic zones: a temperate site (S1), a semi-arid site (S2) and a humid
tropical site (S3). The sites differ in soil texture, reaction and organic
carbon as well as in climate (Table~\ref{tab:sites}). All site and plot data
are synthetic and are generated by a seeded program supplied with the volume.

\input{\subfix{../tables/sites}}

\subsection{Design}
Every trial in this volume uses a randomised complete block design with \res{design-blocks}
blocks at each site \parencite{gomez1984}. Each block contains every treatment
of the trial exactly once, and the order of treatments within each block is
randomised independently. Plots are identified by a code giving the trial,
site, block and position in the block, and every plot of every trial is listed
in a single plot register\index{plot register} (\texttt{data/trials/plot\_register.tsv}), so that
each measured value can be traced to its plot.

\subsection{Analysis at one site}
At a single site the measured response $y_{ij}$ of treatment $i$ in block $j$
is described by the model
\begin{equation}
  y_{ij} = \mu + \tau_i + \beta_j + \varepsilon_{ij},
  \label{eq:rcbd}
\end{equation}
where $\mu$ is the overall mean, $\tau_i$ the effect of treatment $i$,
$\beta_j$ the effect of block $j$ and $\varepsilon_{ij}$ the residual error.
Treatment effects are tested by analysis of variance\index{analysis of variance}, and means are compared
by the least significant difference\index{least significant difference}
\begin{equation}
  \mathrm{LSD} = t_{\alpha/2,\,\nu}\sqrt{\frac{2\,\mathrm{MSE}}{r}},
  \label{eq:lsd}
\end{equation}
where $\mathrm{MSE}$ is the residual mean square with $\nu$ degrees of freedom,
$r$ is the number of blocks and $\alpha = 0.05$.

\subsection{Combined analysis across sites}
When the three sites are analysed together, the model is extended to
\begin{equation}
  y_{ijk} = \mu + \sigma_k + \beta_{j(k)} + \tau_i + (\sigma\tau)_{ik} + \varepsilon_{ijk},
  \label{eq:combined}
\end{equation}
in which $\sigma_k$ is the effect of site $k$, $\beta_{j(k)}$ the effect of
block $j$ within site $k$, and $(\sigma\tau)_{ik}$ the interaction between
treatment and site \parencite{gomez1984}. A significant interaction means that
treatment effects differ between sites; results are then reported site by site
rather than as a single mean.

\subsection{Reporting rules}
Means are reported with the number of decimal places used in the data files.
A treatment effect is described as significant only when the corresponding
$F$-test gives $P < 0.05$, and every table states the number of blocks on
which its means are based.
